% $Id$
% ---------------------------------------------------------------------------
%
% Copyright (c) 2008 Interlegis
%
% Permission is granted to copy, distribute and/or modify this
% document under the terms of the GNU Free Documentation License,
% Version 1.2 or any later version published by the Free Software
% Foundation; with no Invariant Sections, no Front-Cover Texts, and no
% Back-Cover Texts.  A copy of the license is included in the section
% entitled "GNU Free Documentation License". If not, see
% <http://www.gnu.org/licenses/>.
%

\documentclass[a4paper,12pt]{article}
\textheight 22.5cm
\usepackage[brazil]{babel}      % suporte ao portugues brasil
\usepackage[utf8]{inputenc}     % Codificacao dos caracteres (use a
                                % codificacao latin1 caso seu
                                % sistema/editor nao suporte utf8).
\usepackage[T1]{fontenc}
\usepackage{indentfirst}        % indenta o primeiro paragrafo 
\usepackage{graphicx}           % uso para imagens
\usepackage{dsfont}             % fonte

\usepackage{fancyhdr}           % permite alterar o cabecalho
\pagestyle{fancy}               % configura o estilo padrao

\usepackage{verbatim}

% $Id$
% ---------------------------------------------------------------------------
%
%  This is part of the SIGI.
%  Copyright (C) 2008 Interlegis
%  See the file relatorio.tex for copying conditions.
%

\newcommand{\titulo}{Produto III}
\newcommand{\subtitulo}{Sistema de Informações Gerenciais do Interlegis}
\newcommand{\autor}{Guilherme Mesquita Gondim}
\newcommand{\numcontrato}{2008/000471} % Número do contrato caso seja
                                        % consultor.

%
% Local variables:
%   mode: flyspell
%   TeX-master: "relatorio.tex"
% End:
%
                % novos comandos

\usepackage[pdftex,%            % formatacao no PDF
pdfauthor={\autor},%
pdftitle={\titulo \ - \subtitulo},%
]{hyperref}
\hypersetup{backref=true,pdfpagemode=UseOutlines,colorlinks=true,
  breaklinks=true,hyperindex,linkcolor=blue,anchorcolor=black,
  citecolor=blue,filecolor=magenta,menucolor=red,pagecolor=red,
  urlcolor=cyan,bookmarks=true,bookmarksopen=true,pdfpagelayout=SinglePage,
  pdfpagetransition=Dissolve}

\begin{document}
% $Id$
% ---------------------------------------------------------------------------
%
%  This is part of the SIGI.
%  Copyright (C) 2008 Interlegis
%  See the file relatorio.tex for copying conditions.
%

\lhead {
  \setlength{\unitlength}{5mm} 
  \begin{picture}(0,0) 
    \includegraphics[scale=0.65]{imagens/cabecalho.pdf} 
  \end{picture} 
}
\chead{}
\rhead{}
\renewcommand{\headrulewidth}{0pt}

%
% Local variables:
%   mode: flyspell
%   TeX-master: "relatorio.tex"
% End:
%

% $Id$
% ---------------------------------------------------------------------------
%
%  This is part of the SIGI.
%  Copyright (C) 2008 Interlegis
%  See the file relatorio.tex for copying conditions.
%

\textsf{\vspace{6cm}}
\begin{center}
  \noindent
  \large{
    \textbf{\titulo}
  } \\

  \Large{
    \textbf{\subtitulo}
  } \\
  \large{
    \textbf{APO-CASA}
  }

  \vspace{9cm}

  \large{
    \textbf{\autor}\\
    \textsf{Contrato N$^{\circ}$: \numcontrato}\\
  }
\end{center}
\cfoot{}

%
% Local variables:
%   mode: flyspell
%   TeX-master: "relatorio.tex"
% End:
%


% configura o rodape do conteudo pre-textual para numeros romanos
\clearpage
\cfoot{\thepage}
\setcounter{page}{1}
\pagenumbering{Roman}

\tableofcontents
\listoffigures

% configura o rodape das paginas de conteudo para numeros convencionais
\clearpage
\setcounter{page}{1}
\pagenumbering{arabic}

% conteudo
%
%  This is part of the SIGI.
%  Copyright (C) 2008 Interlegis
%  See the file relatorio.tex for copying conditions.
%

\section{Introdução}
Este documento detalha as atividades desenvolvidas durante a segunda
etapa do projeto de desenvolvimento do Sistema de Informações
Gerenciais do Interlegis (SIGI).

A segunda etapa consiste no desenvolvimento do sistema, como projetado
na primeira etapa.

A Seção \ref{sec:info} deste documento possui algumas informações a
respeito do sistema.

A Seção \ref{sec:org} demonstra como o sistema foi organizado,
descreve sua estrutura de diretórios, suas aplicações Django e seus
relacionamentos.

Na Seção \ref{sec:atividades}, as atividades no repositório de código
do sistema são detalhadas.

%
% Local variables:
%   mode: flyspell
%   TeX-master: "relatorio.tex"
% End:
%

%
%  This is part of the SIGI.
%  Copyright (C) 2008 Interlegis
%  See the file relatorio.tex for copying conditions.
%

\section{Desenvolvimento web com Django}
\label{sec:django}

\subsection{Objetivos}
O treinamento visa capacitar desenvolvedores web no uso do framework
Django para a construção de sistemas online.

\subsection{Conteúdo programático}
\begin{enumerate}
\item Introdução
\item Conceitos básicos
  \begin{enumerate}
  \item Model-view-controller (MVC)
  \item Model-template-view (MTV)
  \end{enumerate}
\item Desenvolvimento de uma aplicação
  \begin{enumerate}
  \item Models
  \item Views
  \item Templates
  \end{enumerate}
\item Projetos $x$ Aplicações
  \begin{enumerate}
  \item Organização de projetos
  \item Organização de aplicações
  \end{enumerate}
\end{enumerate}

\subsection{Requisitos}
Para o treinamento, o aluno deverá possuir conhecimentos básicos em
desenvolvimento web e Python.

\subsection{Duração do treinamento}
O treinamento tem duração estimada em 4 horas.

%
% Local variables:
%   mode: flyspell
%   TeX-master: "relatorio.tex"
% End:
%

% $Id$
% ---------------------------------------------------------------------------
%
%  This is part of the SIGI.
%  Copyright (C) 2008 Interlegis
%  See the file relatorio.tex for copying conditions.
%

\section{Requisitos}
\label{sec:requisitos}

\subsection{Detalhamento do Produto}
O produto a ser desenvolvido deverá contemplar as seguintes
funcionalidades:

\begin{itemize}
\item Gerenciamento das informações básicas de identificação das Casas
  Legislativas;
\item Gerenciamento das informações sobre partidos políticos,
  parlamentares, seus mandatos, e legislaturas;
\item Gerenciamento de sessões legislativas, mesas diretora e sua
  composição;
\item Gerenciamento de convênios entre o Interlegis e as Casas
  Legislativas;
\item Gerenciamento dos equipamentos e serviços disponibilizados para
  as Casas Legislativas;
\item Conjunto de autorização e de configuração de perfis de acesso ao
  sistema;
\item Permitir que o usuário efetue pesquisas no banco de dados
  utilizando filtros de seleção;
\item ``Rollback'' (retorno) das alterações, em tempo de
  edição dos dados, nos casos de falhas;
\item Geração de relatórios em PDF, dos dados estratégicos e
  de resultados de pesquisas;
\item Manutenção do histórico das atividades com as Casas Legislativas.
\end{itemize}

Considerar as seguintes características:

\begin{itemize}
\item Interface web compatível com os padrões e recomendações da World
  Wide Web Consortium (W3C);
\item Uso de \emph{AJAX} para validação de dados e para otimização e
  facilidade de uso do sistema;
\item Uso de um \emph{Sistema de Controle de Versões} para gerenciamento e
  acompanhamento da codificação;
\item Ajuda \textit{online} e contextual do sistema;
\item Instalação simplificada;
\item Sistema multi-usuário em rede, concebido para operar em três
  camadas distintas: apresentação (navegador web), aplicação e dados.
\end{itemize}

\subsection{Requisitos de Hardware}
\begin{itemize}
\item Disco: 100 MB ou superior;
\item Memória RAM: 256 MB ou superior;
\item Processador: Pentium III ou superior;
\end{itemize}

\subsection{Requisitos de Software}
\begin{itemize}
\item Sistemas operacionais GNU/Linux ou de família Unix BSD;
\item Servidor web: Apache 2 ou superior;
\item SGBD: PostgreSQL 8.3 ou superior;
\item Linguagem/Interpretador: Python 2.5 ou superior (< 3.0);
\item Framework: Django 0.97-pre (trunk) ou 1.0.
\end{itemize}

\subsection{Requisitos de Performance}
\begin{itemize}
\item Considerar o uso do sistema operacional GNU/Linux ou FreeBSD;
\item Uso do servidor web Apache para a aplicação;
\item Uso do servidor Lighttpd para os arquivos estáticos;
\item Cache das transações com o Banco de Dados e da saída
  (\textit{output}) das requisições HTTP.
\end{itemize}

\subsection{Requisitos de Documentação}
\begin{itemize}
\item Documentação da \textit{Application Programming Interface} (API)
  do software;
\item Manual de configuração e de implantação do sistema;
\item Manual do usuário;
\item Ajuda \textit{online} e contextual do sistema;
\item Arquivos LEIA-ME e de licenciamento.
\end{itemize}

\subsection{Licenças}
O novo SIGI utiliza \emph{softwares livres}, os quais possuem um
licenciamento sem restrições de uso.

Segue abaixo uma lista dos principais \emph{softwares} necessários
para o desenvolvimento do SIGI e suas respectivas licenças:

\begin{description}
\item[Apache:] Apache License 2.0
\item[PostgreSQL:] BSD License
\item[Python:] Python License
\item[Django:] BSD License
\end{description}

O SIGI será publicado como \emph{software livre} utilizando a licença
\textit{GNU General Public License (GNU GPL)} em sua versão 3 e estará
disponível no repositório de \emph{software livre}
Colab\footnote{Endereço web do Colab: http://colab.interlegis.gov.br}
do Interlegis.
%
% Local variables:
%   mode: flyspell
%   TeX-master: "relatorio.tex"
% End:
%

% $Id$
% ---------------------------------------------------------------------------
%
%  This is part of the SIGI.
%  Copyright (C) 2008 Interlegis
%  See the file relatorio.tex for copying conditions.
%

\section{Casos de Uso}
\label{sec:casos}

Esta seção descreve a utilização do sistema através de \emph{Casos de
  Uso}, representando os principais atores e suas interações com o
sistema.

Os Casos de Uso servem de auxílio à compreensão dos usuários quanto à
implementação e uso do sistema.

\subsection{Casos de Uso do SIGI}
A Figura \ref{fig:casos} apresenta os Casos de Uso do SIGI de maneira
simplificada.

\begin{figure}[h]
  \centering
  \includegraphics[width=120mm]{../imagens/casosdeuso.png}
  \caption{Casos de Uso}
  \label{fig:casos}
\end{figure}

\subsubsection{Descrição dos Atores}
\begin{description}
\item[Funcionário Interlegis Autorizado:] usuário apenas com
  atribuições de leitura no sistema. Inclui pessoal da Subsecretaria
  de Apoio Técnico e Relações Institucionais (SSATR) e Subsecretaria
  de Planejamento e Fomento (SSPLF).
\item[SSFAC:] Subsecretaria de Formação e Atendimento à Comunidade do
  Legislativo.
\item[SSADM:] Subsecretaria de Administração.
\item[SEIT:] Serviço de Infra-estrutura Tecnológica.
\end{description}

\subsubsection{Descrição das Atividades}
\begin{description}
\item[Emitir Relatórios:] consiste em obter informações e emitir
  relatórios de diversas partes do sistema.
\item[Emitir Etiquetas:] consiste em obter informações e emitir
  etiquetas de algumas partes do sistema.
\item[Pesquisar e Recuperar Informações:] habilidades de pesquisa e
  obtenção de informações da base de dados do sistema.
\item[Gerenciar Mesas Diretoras:] consiste em inserir, atualizar e
  remover \emph{Mesas Diretoras}, \emph{Sessões Legislativas} e
  modificar a \emph{Composição das Mesas Diretoras}.
\item[Gerenciar Parlamentares:] consiste em inserir, atualizar e
  remover \emph{Parlamentares} e \emph{Partidos}.
\item[Gerenciar Equipamentos:] consiste em inserir, atualizar e
  remover \emph{Equipamentos} e \emph{Fornecedores}.
\item[Gerenciar Casas Legislativas:] consiste em inserir, atualizar e
  remover \emph{Casas Legislativas}.
\item[Gerenciar Convênios:] consiste em inserir, atualizar e remover
  \emph{Convênios}.
\item[Gerenciar Inventário:] consiste em atualizar o \emph{Inventário}
  das Casas Legislativas.
\item[Gerenciar Serviços:] consiste em inserir, atualizar e remover
  \emph{Serviços} prestados às Casas Legislativas.
\end{description}

%
% Local variables:
%   mode: flyspell
%   TeX-master: "relatorio.tex"
% End:
%


\end{document}
